\chapter{Предисловие}

Этот конспект излагает теорию полей и теорию Галуа~--- от расширений полей и многочленов до бесконечной теории
Галуа. Он рассчитан на самостоятельное чтение: все определения даны, все доказательства приведены полностью.

\section*{На основе чего он составлен}

\begin{itemize}[leftmargin=*]
\item \emph{Книга} J.\,S.~Milne, \emph{Fields and Galois Theory}, версия~5.10 (2022). Из неё взяты порядок глав,
нумерация утверждений и основной набор определений и теорем. Номера утверждений в конспекте совпадают с номерами в
книге, поэтому их удобно читать параллельно.
\item \emph{Рукописный конспект} этой книги (269 страниц), который составитель вёл во время чтения. Из него взяты ход
многих рассуждений, рисунки и диаграммы (здесь они перерисованы, кое-где дополнены), замечания на полях. Места, которые
при чтении вызывали вопросы, разобраны здесь подробнее.
\item \emph{Обсуждения прочитанного} с языковой моделью Claude (компания Anthropic): ответы на вопросы, недостающие
шаги доказательств, дополнительные примеры и контрпримеры. Всё существенное из этих обсуждений перенесено в
соответствующие места текста.
\item \emph{Отдельные заметки}, подготовленные по ходу обсуждений: круговое поле $\Q(\zeta_{13})$; группа кватернионов
как группа Галуа; примеры к теории Куммера; норма и след для несепарабельных расширений; примеры бесконечных
расширений Галуа; двойственность Понтрягина и теория Куммера.
\end{itemize}

\section*{Как он сделан}

Текст написан моделью Claude (Claude Opus~5.5) по заданию составителя и под его руководством, глава за главой,
начиная с ноября 2026~года. Составитель определял, что войдёт в конспект и как это должно быть изложено.

\begin{itemize}[leftmargin=*]
\item Утверждения, которых нет в книге, получили номера с буквой (например, 7.8a). Там, где рассуждение источника
неполно или неточно, приведено полное доказательство.
\item Вычисления проверены системой компьютерной алгебры, а каждую главу независимо прочитали два рецензента
(отдельные экземпляры той же модели); найденные ошибки исправлены. Тем не менее опечатки и ошибки возможны.
\item Рисунки сделаны в TikZ. Красным выделено то, что строится в доказательстве, синим~--- исходные данные.
\item Терминология~--- устоявшаяся русская (например, <<проконечная группа>>, <<сепарабельное замыкание>>).
\end{itemize}

\section*{Как читать}

<<Отступления>>~--- пояснения, мотивировки и взгляд вперёд; при первом чтении их можно пропускать. Ссылка вида
<<по 7.4>> отсылает к утверждению 7.4; в электронной версии номера утверждений из уже готовых глав~--- переходы
(они выделены цветом).

Главы добавляются постепенно; сейчас готова глава~7. Ссылки на ещё не написанные главы даны по нумерации книги и
станут переходами, когда эти главы появятся. Исходные тексты (\LaTeX) хранятся вместе с PDF.
